% Web mathematical TeX excerpt reformatted by assistant; not a downloaded native chapter.

\begin{definition}[35.3.1 and 35.3.4]
Let $R\to A$ be a ring map. A descent datum $(N,\varphi)$ for modules with respect to $R\to A$ is given by an $A$-module $N$ and an isomorphism of $A\otimes_R A$-modules
\[
\varphi:N\otimes_R A\longrightarrow A\otimes_R N
\]
such that the cocycle condition holds: the diagram of $A\otimes_R A\otimes_R A$-module maps
\[
\xymatrix{N\otimes_R A\otimes_R A\ar[rr]_{\varphi_{02}}\ar[rd]_{\varphi_{01}}&&A\otimes_R A\otimes_R N\\
&A\otimes_R N\otimes_R A\ar[ru]_{\varphi_{12}}&}
\]
commutes. Here $\varphi_{01}=\varphi\otimes\id_A$, $\varphi_{12}=\id_A\otimes\varphi$, and $\varphi_{02}(n\otimes1\otimes1)=\sum a_i\otimes1\otimes n_i$ if $\varphi(n\otimes1)=\sum a_i\otimes n_i$.

We say a descent datum $(N,\varphi)$ is effective if there exists an $R$-module $M$ and an isomorphism of descent data from $(M\otimes_R A,can)$ to $(N,\varphi)$. The canonical map is
\[
can:(M\otimes_R A)\otimes_R A\longrightarrow A\otimes_R(M\otimes_R A),\quad
(m\otimes a)\otimes a'\longmapsto a\otimes(m\otimes a').
\]
\end{definition}
\paragraph{The identity used from the proof of Lemma 35.3.2.}
Write $p=\sigma^0_0\circ\delta^1_1:N\to N$. By the description of the maps above we deduce that $p$ is also equal to
\[
p=\varphi\otimes\id:N=(N\otimes_R A)\otimes_{(A\otimes_R A)}A
\longrightarrow(A\otimes_R N)\otimes_{(A\otimes_R A)}A=N.
\]
Since $\varphi$ is an isomorphism we see that $p$ is an isomorphism. Write $\varphi(n\otimes1)=\sum\alpha_i\otimes x_i$ for certain $\alpha_i\in A$ and $x_i\in N$. Then $p(n)=\sum\alpha_i x_i$. Next, write $\varphi(x_i\otimes1)=\sum\alpha_{ij}\otimes y_j$ for certain $\alpha_{ij}\in A$ and $y_j\in N$. Then the cocycle condition says that
\[
\sum\alpha_i\otimes\alpha_{ij}\otimes y_j=\sum\alpha_i\otimes1\otimes x_i.
\]
This means that $p(n)=\sum\alpha_i x_i=\sum\alpha_i\alpha_{ij}y_j=\sum\alpha_i p(x_i)=p(p(n))$. Thus $p$ is a projector, and since it is an isomorphism it is the identity.

\begin{lemma}[35.3.7, Tag 039W]
Let $R\to A$ be a faithfully flat ring map. Let $(N,\varphi)$ be a descent datum. Then $(N,\varphi)$ is effective if and only if the canonical map
\[
A\otimes_R H^0(s(N_\bullet))\longrightarrow N
\]
is an isomorphism.
\end{lemma}
\begin{proof}
If $(N,\varphi)$ is effective, then we may write $N=A\otimes_R M$ with $\varphi=can$. It follows that $H^0(s(N_\bullet))=M$ by Lemmas 35.3.3 and 35.3.6. Conversely, suppose the map of the lemma is an isomorphism. In this case set $M=H^0(s(N_\bullet))$. This is an $R$-submodule of $N$, namely
\[
M=\{n\in N\mid1\otimes n=\varphi(n\otimes1)\}.
\]
The only thing to check is that via the isomorphism $A\otimes_R M\to N$ the canonical descent data agrees with $\varphi$. We omit the verification.
\end{proof}
\begin{lemma}[35.3.8, Tag 039X]
Let $R\to A$ be a faithfully flat ring map, and let $R\to R'$ be faithfully flat. Set $A'=R'\otimes_R A$. If all descent data for $R'\to A'$ are effective, then so are all descent data for $R\to A$.
\end{lemma}
\begin{proof}
Let $(N,\varphi)$ be a descent datum for $R\to A$. Set $N'=R'\otimes_R N=A'\otimes_A N$, and denote $\varphi'=\id_{R'}\otimes\varphi$ the base change of the descent datum $\varphi$. Then $(N',\varphi')$ is a descent datum for $R'\to A'$ and $H^0(s(N'_\bullet))=R'\otimes_R H^0(s(N_\bullet))$. Moreover, the map $A'\otimes_{R'}H^0(s(N'_\bullet))\to N'$ is identified with the base change of the $A$-module map $A\otimes_R H^0(s(N))\to N$ via the faithfully flat map $A\to A'$. Hence we conclude by Lemma 35.3.7.
\end{proof}
\begin{proposition}[35.3.9, part (1), Tag 023N]
Let $R\to A$ be a faithfully flat ring map. Then any descent datum on modules with respect to $R\to A$ is effective.
\end{proposition}
\begin{proof}
We only prove (1) and omit the proofs of (2) and (3). As $R\to A$ is faithfully flat, there exists a faithfully flat base change $R\to R'$ such that $R'\to A'=R'\otimes_R A$ has a section (namely take $R'=A$ as in the proof of Lemma 35.3.6). Hence, using Lemma 35.3.8 we may assume that $R\to A$ has a section, say $\sigma:A\to R$. Let $(N,\varphi)$ be a descent datum relative to $R\to A$. Set
\[
M=H^0(s(N_\bullet))=\{n\in N\mid1\otimes n=\varphi(n\otimes1)\}\subset N.
\]
By Lemma 35.3.7 it suffices to show that $A\otimes_R M\to N$ is an isomorphism.

Take an element $n\in N$. Write $\varphi(n\otimes1)=\sum a_i\otimes x_i$ for certain $a_i\in A$ and $x_i\in N$. By Lemma 35.3.2 we have $n=\sum a_i x_i$ in $N$ (because $\sigma^0_0\circ\delta^1_1=\id$ in any cosimplicial object). Next, write $\varphi(x_i\otimes1)=\sum a_{ij}\otimes y_j$ for certain $a_{ij}\in A$ and $y_j\in N$. The cocycle condition means that
\[
\sum a_i\otimes a_{ij}\otimes y_j=\sum a_i\otimes1\otimes x_i
\]
in $A\otimes_R A\otimes_R N$. We conclude two things from this:
\begin{enumerate}
\item applying $\sigma$ to the first $A$ we get $\sum\sigma(a_i)\varphi(x_i\otimes1)=\sum\sigma(a_i)\otimes x_i$,
\item applying $\sigma$ to the middle $A$ we get $\sum_i a_i\otimes\sum_j\sigma(a_{ij})y_j=\sum a_i\otimes x_i$.
\end{enumerate}
Part (1) shows that $\sum\sigma(a_i)x_i\in M$. Applying this to $x_i$ we see that $\sum\sigma(a_{ij})y_i\in M$ for all $i$. Multiplying out the equation in (2) we conclude that $\sum_i a_i(\sum_j\sigma(a_{ij})y_j)=\sum a_i x_i=n$. Hence $A\otimes_R M\to N$ is surjective. Finally, suppose that $m_i\in M$ and $\sum a_i m_i=0$. Then we see by applying $\varphi$ to $\sum a_i m_i\otimes1$ that $\sum a_i\otimes m_i=0$. In other words $A\otimes_R M\to N$ is injective and we win.
\end{proof}
